\begin{center}
\begin{tikzpicture}[x=.88cm,y=.88cm,font=\small]
  \fill[paper] (0,0) rectangle (12.3,6.6);

  \draw[source!70,line width=.8pt] (.75,2.2) .. controls (3.4,2.58) and
    (7.05,1.8) .. (11.72,2.27);
  \node[text=source,below,font=\scriptsize] at (8.1,2.06) {黄河方向（示意）};

  \fill[dynasty!43] (7.1,4.65) ellipse (1.65 and 1.12);
  \node[align=center,text=white] at (7.1,4.65) {北魏\\平城为中心\\的扩张网络};

  \fill[timeline!25] (4.0,3.35) ellipse (1.22 and .83);
  \node[align=center] at (4.0,3.35) {夏\\统万城};

  \fill[source!26] (2.0,4.92) ellipse (1.2 and .83);
  \node[align=center] at (2.0,4.92) {北凉\\姑臧、河西};

  \fill[timeline!24] (10.5,4.85) ellipse (1.16 and .84);
  \node[align=center] at (10.5,4.85) {北燕\\龙城};

  \fill[source!20] (8.53,2.88) ellipse (1.05 and .62);
  \node[align=center,font=\scriptsize] at (8.53,2.88) {河北、辽西\\旧战区};

  \draw[->,very thick,color=dynasty]
    (6.0,4.18) .. controls (5.37,3.85) and (4.91,3.59) .. (4.72,3.48);
  \node[text=dynasty,below,font=\scriptsize] at (5.24,3.75)
    {427--431：取统万城，夏亡};

  \draw[->,very thick,color=dynasty]
    (8.62,4.71) .. controls (9.05,4.8) and (9.35,4.85) .. (9.38,4.85);
  \node[text=dynasty,above,font=\scriptsize] at (8.87,5.24)
    {436：北燕亡};

  \draw[->,very thick,color=dynasty]
    (5.6,4.93) .. controls (4.54,5.05) and (3.73,5.0) .. (3.23,4.97);
  \node[text=dynasty,above,font=\scriptsize] at (4.37,5.35)
    {439：北凉亡};

  \draw[->,thick,dashed,color=timeline]
    (7.64,3.7) .. controls (7.93,3.45) and (8.06,3.3) .. (8.12,3.24);
  \node[text=timeline,right,font=\scriptsize,align=left] at (8.92,3.5)
    {先后夺取\\并非一条连续前线};

  \node[draw=source!75,fill=paper,rounded corners=2pt,align=center,
    inner sep=3pt,font=\scriptsize] at (5.8,.86)
    {427--439年的北魏扩张，终结了北方若干主要割据中心；\\
    “统一”不等于地方控制即时、均质地完成};
\end{tikzpicture}

\smallskip
\small\textit{图\quad 北魏完成对北方若干主要割据中心的征服，约427--439年：据《魏书·世祖纪上、下》、〈赫连昌传〉、〈沮渠蒙逊传〉、〈冯弘传〉及《资治通鉴》卷一百二十至一百二十三，概括统万城、龙城与姑臧的先后失守。色块不绘制北魏或各国的实测疆界；城市定位、河流和征服箭线均为约略示意，且不表示全部行军路线、驻军范围或征服后行政控制的精确速度。}
\end{center}

\phantomsection\label{fig:northern-wei-unification}
